\chapter{The Hunt}\label{ch:B2}

This chapter answers \textbf{B2} of Part B (17 marks): finding the planted anomalies in our
captures, saying how many there were, and giving a display filter for each that returns exactly
the frames we claim.

\textbf{Reference window: D2\_W2} (21-08-2026, 14:55--15:08 IST). We use one fixed window for
both anomaly tables because it is the window in which both teammates hold complete, non-zero
captures taken three minutes apart, which makes frame-level comparison between the two devices
meaningful. Every filter below was run in \texttt{tshark} against the submitted
\texttt{pcapng} files; the ``Frames returned'' column is the count that filter actually produces
on that file, not an estimate.

\section{B2.1: How many were planted?}

The handout says between three and six, and that working out how many is part of the job. We did
not assume the two seeds matched, and they do not: \textbf{2024CS10300 has four, 2024CS10582 has
six.}

We tested all six flavours listed in the handout against both captures and recorded the
negatives as well as the positives, because ``not in this seed'' is a result.

\begin{table}[H]
\centering
\caption{The six handout flavours tested against both captures. Two are genuinely absent from
2024CS10300's seed; the evidence for those absences is given below.}
\label{tab:b2flavours}
\footnotesize
\begin{tabular}{@{}>{\raggedright\arraybackslash}p{5.6cm}>{\raggedright\arraybackslash}p{3.6cm}>{\raggedright\arraybackslash}p{3.6cm}@{}}
\toprule
Flavour (handout wording) & 2024CS10300 & 2024CS10582 \\
\midrule
Something repeating on a fixed schedule & found: UDP 20020, 11\,s & found: UDP 20023, 9\,s \\
A hostname far longer and stranger than any real one & found: token NVAAAR56 & found: token 2GEYWNDD \\
A password sitting in the clear & found: \texttt{mfarooq24} & found: \texttt{avinash36} \\
A rapid burst of tiny packets across many port numbers & found: 24 ports & found: 48 ports \\
A packet whose TTL is nothing like its neighbours & \textbf{absent} & found: TTL 17 \\
Traffic using one protocol's port number for something else & \textbf{absent} & found: HTTP on 10443 \\
\bottomrule
\end{tabular}
\end{table}

\subsection{Evidence for the two absences on 2024CS10300}

\textbf{Odd TTL: absent.} Every packet this host sent to 168.144.154.92 carries
\texttt{ip.ttl == 128}, the Windows default. The filter
\texttt{ip.dst == 168.144.154.92 \&\& ip.ttl != 128 \&\& !icmp} returns \textbf{0 frames} in all
five windows. The \texttt{!icmp} term is doing real work: without it the same filter returns 10,
because Wireshark also matches the IP header quoted inside the server's ICMP port-unreachable
replies, where the TTL has decremented to 116 in transit. That is an artefact of how the filter
is written, not a planted anomaly.

\textbf{Wrong-port protocol: absent.} The complete set of destination ports this host used
toward the netlab server is 5763--5786, 8080 and 20020. Port 8080 is a conventional HTTP
alternate port carrying HTTP, so nothing here is a protocol/port mismatch. Contrast anomaly B6
below, where plaintext HTTP appears on 10443.

\section{B2.2: Anomaly table, 2024CS10300 (SEED\_TAG 5A11D1FB)}

\begin{table}[H]
\centering
\caption{Four planted anomalies on 2024CS10300. All rows refer to
\texttt{2024CS10300\_D2\_W2.pcapng}.}
\label{tab:b2anom10300}
\scriptsize
\begin{tabular}{@{}c>{\raggedright\arraybackslash}p{3.0cm}c>{\raggedright\arraybackslash}p{2.9cm}>{\raggedright\arraybackslash}p{3.5cm}c@{}}
\toprule
\# & What it is & Frame & Value it encodes & Display filter & Frames \\
\midrule
A1 & UDP beacon to the netlab server on a fixed schedule, 95-byte frames, each answered with ICMP port-unreachable & 8543 & \textbf{Interval = 11\,s} (10 sends; deltas 10.907--10.946\,s) & \texttt{udp.dstport == 20020 \&\& ip.dst == 168.144.154.92 \&\& !icmp} & 10 \\
\addlinespace
A2 & DNS A queries for eight sequential nonsense hostnames under a lab subdomain & 8603 & \textbf{TOKEN = NVAAAR56}, base32 of the label \texttt{jzlecqkbki2tm} & \texttt{dns.flags.response == 0 \&\& dns.qry.name contains \qq .probe.col334.lab\qq} & 8 \\
\addlinespace
A3 & HTTP POST carrying credentials in the clear over port 8080 & 8975 & \texttt{username=mfarooq24} \texttt{password=zermxjpc6g} \texttt{remember=1} & \texttt{http.request.method == \qq POST\qq\ \&\& http.request.uri == \qq /portal/login\qq} & 1 \\
\addlinespace
A4 & Burst of 24 tiny SYNs to consecutive ports, $\sim$313\,ms apart, every one answered RST/ACK & 8628 & Ports \textbf{5763--5786}, 66-byte SYNs, no port open & \texttt{tcp.flags.syn == 1 \&\& tcp.flags.ack == 0 \&\& ip.dst == 168.144.154.92 \&\& tcp.dstport $>$= 5763 \&\& tcp.dstport $<$= 5786} & 24 \\
\bottomrule
\end{tabular}
\end{table}

\section{B2.3: Anomaly table, 2024CS10582 (SEED\_TAG 5E857C61)}

\begin{table}[H]
\centering
\caption{Six planted anomalies on 2024CS10582. All rows refer to
\texttt{2024CS10582\_D2\_W2.pcapng}.}
\label{tab:b2anom10582}
\scriptsize
\begin{tabular}{@{}c>{\raggedright\arraybackslash}p{3.0cm}c>{\raggedright\arraybackslash}p{2.9cm}>{\raggedright\arraybackslash}p{3.5cm}c@{}}
\toprule
\# & What it is & Frame & Value it encodes & Display filter & Frames \\
\midrule
B1 & UDP beacon on a fixed schedule, 113-byte frames & 3506 & \textbf{Interval = 9\,s} (10 sends; deltas 9.000\,s) & \texttt{udp.dstport == 20023 \&\& ip.dst == 168.144.154.92 \&\& !icmp} & 10 \\
\addlinespace
B2 & DNS queries for eight nonsense hostnames, asked A and AAAA, bare and with the \texttt{iitd.ac.in} search suffix & 4249 & \textbf{TOKEN = 2GEYWNDD}, base32 of the label \texttt{gjdukwkxjzcei} & \texttt{dns.flags.response == 0 \&\& dns.qry.name contains \qq .probe.col334.lab\qq} & 32 \\
\addlinespace
B3 & HTTP POST carrying credentials in the clear over port 8080 & 4600 & \texttt{username=avinash36} \texttt{password=bdrhnygavp} \texttt{remember=1} & \texttt{http.request.method == \qq POST\qq\ \&\& http.request.uri == \qq /portal/login\qq} & 1 \\
\addlinespace
B4 & Burst of 48 tiny SYNs to consecutive ports, $\sim$50\,ms apart & 4401 & Ports \textbf{5370--5417}, 74-byte SYNs & \texttt{tcp.flags.syn == 1 \&\& tcp.flags.ack == 0 \&\& ip.dst == 168.144.154.92 \&\& tcp.dstport $>$= 5370 \&\& tcp.dstport $<$= 5417} & 48 \\
\addlinespace
B5 & Seven 51-byte UDP datagrams sent with TTL 17 while every other packet this host sends leaves at TTL 64 & 4564 & \textbf{TTL = 17}, dst port 30021, $\sim$0.8\,s apart & \texttt{ip.dst == 168.144.154.92 \&\& ip.ttl == 17 \&\& udp.dstport == 30021} & 7 \\
\addlinespace
B6 & Plaintext HTTP \texttt{GET /status} on TCP 10443, a port commonly used for HTTPS-alt services & 4609 & Cleartext HTTP on an HTTPS-alt port; server replies \texttt{Apache/2.4.58 (Ubuntu)} & \texttt{tcp.dstport == 10443 \&\& http.request.method == \qq GET\qq} & 5 \\
\bottomrule
\end{tabular}
\end{table}

\section{B2.4: Why these filters and not looser ones}

\textbf{A1 / B1: the \texttt{!icmp} term is necessary here.} Each beacon datagram is answered
by an ICMP port-unreachable that quotes the original UDP header, and Wireshark matches
\texttt{udp.dstport} inside that quote. Without \texttt{!icmp} the filter returns 20 rather than
10, and the count stops meaning ``beacons sent''. Verified on
\texttt{2024CS10300\_D2\_W2}: 20 without, 10 with. We use \texttt{ip.dst ==} rather than
\texttt{ip.addr ==} for the same reason: direction matters to what the count means.

\textbf{A4 / B4: the \texttt{!icmp} term is \emph{not} necessary here, and we checked.} The
obvious worry is that the beacon problem repeats for the SYN burst. It does not, because the
netlab server answers these SYNs at layer 4, not with an ICMP error: frame 8629 in
\texttt{2024CS10300\_D2\_W2.pcapng} is TCP 5763 $\rightarrow$ 61994 with the RST and ACK flags
set. No ICMP is generated anywhere in the burst, so there is no quoted TCP header to contaminate
the count. Running both forms confirms it: 24 with and without \texttt{!icmp} for A4, 48 with and
without for B4. We left the redundant term out.

What the filter \emph{does} need is the port range. \texttt{ip.dst == 168.144.154.92} alone would
sweep in the beacon, the POST and (on 2024CS10582) the 10443 traffic. Constraining to
SYN-without-ACK plus the observed range isolates the burst exactly and makes the returned count
equal the number of ports probed.

\textbf{B5: scoping the TTL, and a detail worth knowing.} \texttt{ip.ttl == 17} alone is not
specific: the traceroutes in the same capture emit probes at every TTL from 1 to 20, so hop-17
probes match too. Adding \texttt{ip.dst == 168.144.154.92 \&\& udp.dstport == 30021} separates
the planted flow from our own measurement traffic. This is a case where the obvious filter from
the handout's starter list gives a correct-looking but contaminated answer.

\texttt{!icmp} is again unnecessary, and the reason is instructive. The ICMP port-unreachables do
exist (frames 4565, 4567 and 4569 answer frames 4564, 4566 and 4568), but the IP header
quoted inside them carries \texttt{ip.ttl == 5}, not 17. That is the TTL as the datagram
\emph{arrived at the server}: sent at 17, arrived at 5, so it crossed 12 hops and still reached
the destination with margin to spare. So \texttt{ip.ttl == 17} cannot match the quote, and the
filter returns 7 either way. The anomaly is a deliberately low initial TTL, not a packet dying in
transit.

\textbf{A2 / B2: the handout's starter filter misses ours.} \texttt{dns.qry.name.len $>$ 40}
returns \textbf{0} planted frames on 2024CS10300. Windows queried the bare name
\texttt{jzlecqkbki2tm-00.probe.col334.lab}, 33 characters, because the host did not append its
\texttt{iitd.ac.in} DNS suffix to these lookups. 2024CS10582's Ubuntu host did append it,
producing 44-character names the starter filter catches. The anomaly is not invisible on Windows
(our filter finds all eight), but the \emph{suggested length filter} behaves differently on
the two machines purely because of resolver configuration. Had we trusted it we would have
reported three anomalies on this device and never recovered the token.

We match on \texttt{.probe.col334.lab} rather than on the server IP or on \texttt{probe}, so the
filter picks up the planted lookups and nothing else, and \texttt{dns.flags.response == 0} drops
the replies so the count equals queries emitted.

\textbf{B6: a note on dissector dependence.} Port 10443 is not in Wireshark's default HTTP
port list, so \texttt{http.request.method} resolves here only because Wireshark's HTTP heuristic
dissector is enabled by default. If it is disabled, the filter parses but returns 0. A
dissector-independent equivalent, verified to return the same 5 frames, is:

\begin{quote}\ttfamily\footnotesize
\begin{verbatim}
tcp.dstport == 10443 && tcp.payload[0:3] == 47:45:54
\end{verbatim}
\end{quote}

\noindent(\texttt{47 45 54} is \texttt{GET} in ASCII.) We give the readable filter as primary and
this as a fallback.

\section{B2.5: Frames returned is a property of the capture, not of the plant}

The planted signatures recur across every window we hold, but the frame counts are not identical
everywhere, and the one exception is worth reporting rather than smoothing over.

Filter A2 returns \textbf{10 frames in 2024CS10300\_D1\_W2} and 8 in every other window. The
extra two are not extra plants:

\begin{table}[H]
\centering
\caption{Four frames, two lookups: the Windows resolver failing over to the secondary server.}
\label{tab:b2failover}
\footnotesize
\begin{tabular}{@{}ccll>{\raggedright\arraybackslash}p{5.6cm}@{}}
\toprule
Frame & t (s) & Sent to & DNS txn ID & Query \\
\midrule
1375 & 422.859 & 10.7.172.202 & 0xda45 & \texttt{jzlecqkbki2tm-03.probe.col334.lab} \\
1376 & 422.884 & 10.10.1.4 & 0xda45 & \texttt{jzlecqkbki2tm-03.probe.col334.lab} \\
1383 & 427.440 & 10.7.172.202 & 0x3dc1 & \texttt{jzlecqkbki2tm-06.probe.col334.lab} \\
1384 & 427.465 & 10.10.1.4 & 0x3dc1 & \texttt{jzlecqkbki2tm-06.probe.col334.lab} \\
\bottomrule
\end{tabular}
\end{table}

Same transaction IDs, 25\,ms apart, re-sent to the \textbf{secondary} resolver. The primary
(10.7.172.202) did not answer quickly enough and the Windows resolver failed over to 10.10.1.4.
Eight planted lookups, ten query packets.

This is the second time this device's resolver has changed what the anomaly looks like on the
wire, the first being the missing search suffix in Section~B2.4. Neither difference is in the
plant; both are in the host. It is a reminder that a frame count is a measurement of our own
machine as much as of the thing we are hunting.

\section{B2.6: Required per-teammate block}

\begin{quote}\ttfamily\footnotesize
\begin{verbatim}
ROLL:     2024CS10300
SEED_TAG: 5A11D1FB
TOKEN:    NVAAAR56
INTERVAL: 11
\end{verbatim}
\end{quote}

\begin{quote}\ttfamily\footnotesize
\begin{verbatim}
ROLL:     2024CS10582
SEED_TAG: 5E857C61
TOKEN:    2GEYWNDD
INTERVAL: 9
\end{verbatim}
\end{quote}

\textbf{How the tokens decode.} The hostname label is base32 with the padding stripped and the
case lowered. Restore it: upper-case the label, append \texttt{=} until the length is a multiple
of 8, then base32-decode. All eight planted queries carry the same encoded label, so the token
can be recovered from any one of them; we verified the label is byte-identical across all eight
probes and across every window.

\begin{quote}\ttfamily\footnotesize
\begin{verbatim}
import base64
base64.b32decode("JZLECQKBKI2TM" + "===")   # -> b'NVAAAR56'   (2024CS10300)
base64.b32decode("GJDUKWKXJZCEI" + "===")   # -> b'2GEYWNDD'   (2024CS10582)
\end{verbatim}
\end{quote}

\section{B2.7: Cross-device observation}

The two seeds share four flavours: periodic UDP beacon, long DNS hostname, cleartext HTTP POST
and rapid SYN burst. The TTL anomaly and the wrong-port HTTP service occur only on
2024CS10582. Where the four shared flavours overlap, the seeded parameters differ between the two
rolls: beacon port and interval, token, credentials, and the position, width and spacing of the
port burst.

Running both tables side by side is what let us be confident that ``absent'' on 2024CS10300 means
absent rather than missed. We knew exactly what a TTL anomaly and a wrong-port service look like
in this dataset, because the other device had one of each, so we could search for that specific
signature and record its absence with a filter and a count rather than a shrug.

\section{B2.8: Confirmation in the re-recorded D2W3}

Both teammates re-recorded D2W3 on 22-08 (2024CS10300 20:57:38--21:07:38 IST, 2024CS10582
20:58:57--21:08:58 IST, 79 seconds apart, both a full 600.2\,s by the pcapng Interface Statistics
Block). Every filter in the two tables above was re-run against those files and returns exactly
the counts stated:

\begin{table}[H]
\centering
\caption{Every filter re-run against the re-recorded D2W3 captures. Counts match the reference
window exactly.}
\label{tab:b2d2w3}
\footnotesize
\begin{tabular}{@{}lccccccc@{}}
\toprule
File & A1/B1 & A2/B2 & A3/B3 & A4/B4 & B5 & B6 & negative TTL check \\
\midrule
\texttt{2024CS10300\_D2\_W3} & 10 & 8 & 1 & 24 & --- & --- & 0 \\
\texttt{2024CS10582\_D2\_W3} & 10 & 32 & 1 & 48 & 7 & 5 & --- \\
\bottomrule
\end{tabular}
\end{table}

The D1W2 resolver failover described in Section~B2.5 does not recur here: A2 returns 8, the plant
count, so the primary resolver answered every query in this window.
